\documentclass[11pt]{article}
% Shared preamble for the preseason 2026 articles. House style follows papers/queta.tex.
\usepackage[margin=1in]{geometry}
\usepackage{amsmath,amssymb}
\usepackage{booktabs}
\usepackage{graphicx}
\usepackage{xcolor}
\usepackage{multirow}
\usepackage{pgfplots}
\pgfplotsset{compat=1.17}
\usepgfplotslibrary{groupplots}
\usepackage{caption}
\captionsetup{font=small,labelfont=bf}
\usepackage[colorlinks=true,linkcolor=black,citecolor=black,urlcolor=blue!50!black]{hyperref}
\usepackage{microtype}
\usepackage{xurl}
\usepackage{enumitem}
\usepackage{float}

\definecolor{court}{HTML}{0E5A3A}
\definecolor{amberx}{HTML}{B26A1F}
\definecolor{steel}{HTML}{5E6E75}
\definecolor{fade}{HTML}{AEB3B5}

\newcommand{\gravity}{\textsc{gravity}}
\newcommand{\ctg}{\textsc{ctg}}
\newcommand{\plainbox}[1]{\par\medskip\noindent\begin{center}%
\colorbox{court!8}{\parbox{0.93\textwidth}{\small \textbf{\textcolor{court}{In plain terms:}} #1}}%
\end{center}\medskip}

% Inline source marker: \src{n}{url} prints a linked superscript [n].
\newcommand{\src}[2]{\textsuperscript{\href{#2}{[#1]}}}

\usepackage{fancyhdr}
\newcommand{\statusfoot}{%
  \fancyhf{}%
  \fancyfoot[L]{\footnotesize\textcolor{steel}{An article, not a research paper}}%
  \fancyfoot[C]{\thepage}%
  \fancyfoot[R]{\footnotesize\textcolor{steel}{thegravityreport.com}}%
  \renewcommand{\headrulewidth}{0pt}%
  \renewcommand{\footrulewidth}{0pt}%
}
\fancypagestyle{gravity}{\statusfoot}
\fancypagestyle{plain}{\statusfoot}
\pagestyle{gravity}

% Version line required at the top of every preseason article.
\newcommand{\versionline}[1]{\par\noindent\begin{center}%
\colorbox{fade!25}{\parbox{0.93\textwidth}{\centering\small
\textbf{v1 --- methodology and figures verified against public sources as of #1.
This article will be revised.}\\[2pt]
Not a research paper: an article, typeset in \LaTeX. Independent analysis, not peer reviewed.}}%
\end{center}\smallskip}

\newcommand{\betbox}[1]{\par\medskip\noindent\begin{center}%
\fcolorbox{court}{court!6}{\parbox{0.93\textwidth}{\small #1}}%
\end{center}\medskip}

\setlist{nosep,leftmargin=1.4em}
\setlength{\emergencystretch}{3em}

% Forecast display: plain-text claim, a bar filled to the probability, a tick at 50% (coin flip).
% Usage: \forecastitem{50}{Claim text.}{Grading note.}
\newcommand{\forecastscale}{%
  \par\noindent\begin{tikzpicture}
  \node[font=\scriptsize, text=steel, anchor=west] at (0,0) {0\%};
  \node[font=\scriptsize, text=steel] at (4.5,0) {coin flip};
  \node[font=\scriptsize, text=steel, anchor=east] at (9,0) {100\%};
  \end{tikzpicture}\par\vspace{-2pt}}
\newcommand{\forecastitem}[3]{%
  \par\medskip\noindent{\small\textbf{#2}}\par\vspace{3pt}\noindent
  \begin{tikzpicture}[baseline=(b.base)]
  \fill[fade!30] (0,0) rectangle (9,0.34);
  \fill[court!80] (0,0) rectangle ({9*#1/100},0.34);
  \draw[steel, line width=0.6pt] (4.5,-0.07) -- (4.5,0.41);
  \node[anchor=west, font=\large\bfseries, text=court] (b) at (9.25,0.17) {#1\%};
  \end{tikzpicture}%
  \ifx\relax#3\relax\else\par\vspace{1pt}\noindent{\footnotesize\textcolor{steel}{#3}}\fi\par}


\title{\textbf{The Kawhi Anomaly}\\[4pt]
\large The biggest rebound ever recorded by a 34-year-old\\[2pt]
\normalsize A Gravity Report Article}
\author{}
\date{October 6, 2026}

\begin{document}
\maketitle
\vspace{-2.5em}
\versionline{October 6, 2026}

\begin{abstract}
\noindent
In the 2024--25 season, Kawhi Leonard played 37 games and was graded as a good starter. In
the 2025--26 season, at 34 years old, he played 65 games and was graded as the 5th best player
in the league per minute. His jump in impact score (+1.31) is the biggest one-year gain ever
recorded by a player 34 or older, and he is the only star since 1978 to set career highs in
usage and true shooting in the same season after 33. Late surges have an unkind history. Only
eight players have surged to star level after 31. None of them kept $\tfrac{3}{4}$ of that gain,
and the median kept almost none of it. Leonard's surge came in two parts, taking on a heavier
role and hotter shooting. Both of those parts age differently historically. When we rebuild his
season piece by piece, we predict an impact of +1.31 next season, roughly a top-ten level, which
means he keeps about a third of his surge.
\end{abstract}

\section*{Two records}

The latest edition of The 644 ranks Leonard at sixth with a score of 77.3, the oldest player in
the top ten by three years.\src{1}{https://thegravityreport.com/the644/the644-2026-07.csv} The
season that earned him that score came out of nowhere. We ran every NBA season since 1977--78
through the model's public variant, \gravity-RS. We use this version because the full 644 leans
on on/off data from Cleaning the Glass, a subscription stats site that strips out garbage time
(minutes when the game is already decided). That data only goes back to 2003--04 and can't be
shared.
\gravity-RS, short for \gravity{} Retrospective, is the same formula without that one input, so
it works for every season and anyone can check it. It reproduces the historical values in
Articles No.~1 and No.~2 to within 0.03.\src{2}{https://www.basketball-reference.com/leagues/}
Two records fall out of that table.

Before going further, a quick note on what the numbers mean. Every rating in this article is an
impact score: how much a player helps his team per minute, compared with an average NBA minute.
Zero is average. A plus number is better than average, and the bigger it is, the better. To give
you a sense of the scale from last season: +0.6 was better than 80\% of the league's regular
rotation players, +1.3 was roughly a top-ten player, and +2.0 or more was a top-six player. Only
five players got there.

The first record is the size of the jump. Leonard went from +0.86, which was better than 84\% of
rotation players, to +2.17, the fifth-best in the league, a gain of 1.31. Among players 34 or
older with 1,500 minutes, nobody has ever improved their impact score by more in one year. P.J.
Tucker is the next closest, at +1.27, and he climbed from below average to just average. Only one
player aged 32 or older has made a bigger jump, Damian Lillard in 2022--23, and he gave most of it
back the next season, keeping only 17\% of the improvement.

The second is the shape of the season. Leonard used 33.5\% of the Clippers' possessions and shot
.629 true shooting, and both numbers were career
highs.\src{3}{https://www.basketball-reference.com/players/l/leonaka01.html} Since 1978, four
players have set career highs in these categories in the same season after turning 33. The other
three players are Bojan Bogdanović (2022--23), Elden Campbell (2001--02), and John Lucas
(1986--87). However, none of them were graded higher than +1.04 that year. Kawhi was graded
+2.17.

\section*{What happens to late surges}

A 34-year-old who suddenly plays better than he did at 33 raises an obvious question: is he
reaching a new level, or is it just a good year? Historically, we have an answer for the general
case. Since 1978, 70 players aged 31 or older have improved their impact score by 0.6 or more in a
single season, about the distance between an average player and one better than 80\% of the
league. For the 66 of those who played the next year, the median player dropped by 0.57. The
median player kept a third of his gain, and 13 of them kept three-quarters or more. Those 13 are
mostly role players who climbed towards average; however, the stars tell a much harsher story.

\begin{table}[H]
\centering\small
\caption{Every player since 1977--78 who improved by 0.6 or more after turning 31 and reached a
star-level +1.5. ``Kept'' is the share of that one-season improvement still there the next
season.}
\begin{tabular}{llccccr}
\toprule
Player & Season & Age & Before & Surge & Next year & Kept \\
\midrule
Clyde Drexler & 1994--95 & 32 & $+0.81$ & $+1.81$ & $+1.07$ & 26\% \\
Chauncey Billups & 2007--08 & 31 & $+1.51$ & $+2.21$ & $+1.07$ & $-62$\% \\
Kevin Garnett & 2007--08 & 31 & $+1.51$ & $+2.27$ & $+1.50$ & $-2$\% \\
Steve Nash & 2009--10 & 35 & $+0.85$ & $+1.51$ & $+1.15$ & 45\% \\
Tim Duncan & 2012--13 & 36 & $+1.00$ & $+1.60$ & $+0.98$ & $-3$\% \\
Dirk Nowitzki & 2013--14 & 35 & $+0.89$ & $+1.57$ & $+0.70$ & $-28$\% \\
Jimmy Butler & 2022--23 & 33 & $+1.82$ & $+2.47$ & $+1.35$ & $-71$\% \\
Damian Lillard & 2022--23 & 32 & $+0.45$ & $+2.15$ & $+0.73$ & 17\% \\
\midrule
\textbf{Kawhi Leonard} & 2025--26 & 34 & $+0.86$ & $\mathbf{+2.17}$ & ? & \\
\bottomrule
\end{tabular}
\label{tab:stars}
\end{table}

Eight players surged to star level after 31. Not one of them kept three-quarters of the
improvement they made in their surge season the following year. Nash kept the most with 45\%.
Only five finished the next season at or below where they started. This record alone would
settle the question of whether Leonard's surge would stay next year, but it doesn't. Kawhi's
surge is entirely different.

\section*{The anatomy of the surge}

\gravity-RS builds a season from nine measured parts, and some of them repeat from one year to the
next more reliably than others. Think of it as a reliability score from 0 to 1. At 1, last season
tells you exactly what to expect next season. At 0, it tells you nothing. Among players 31 and
older, a player's scoring volume, creation, and offensive load all score above 0.9, and shooting
efficiency scores 0.72, the least reliable part of the whole thing.


This table splits Leonard's surge into what he did more of, and what he did better. Last year, he
took on far more of the offense. His load and volume each rose by about a full point on that
scale, and those are the parts that usually hold. His efficiency also jumped from about average
to +1.5, close to elite, and that's the part that's supposed to hold least. The two box-score
numbers the model leans on most, Box Plus/Minus and win shares, sit in between: more reliable than
shooting, less reliable than role.

Cleaning The Glass, which filters out garbage time, shows where the efficiency came from. Kawhi
scored 126.9 points per 100 shot attempts, the best mark of his career and putting him in the 91st
percentile at his position. Half of his shots came from mid-range, as they have for a decade, and
he made 52\% of them, also a career high. He drew shooting fouls on 13.9\% of his shots, up from
8.7\% a year earlier. However, his three-point percentage fell from 41\% to
38\%.\src{4}{https://cleaningtheglass.com/stats/player/2173}

That shows us two different kinds of improvements. The foul drawing might glimpse a return to his
own prime. He drew fouls on 13\% to 15\% of his shots every season from 2016 to 2020, and getting
to the free-throw line is one of the more reliable parts of an older player's game, scoring 0.86
on that same 0-to-1 scale. The mid-range makes are a career high on the shot he takes most;
however, the shooting efficiency is the part that carries over least. If the forecast we make is
wrong, the mid-range shot will most likely be where it goes wrong.

Put simply, our forecast expects him to keep about half to 60\% of his bigger role in the offense but
only about a third of his shooting jump. We predict a 2026--27 impact of +1.31, roughly a top-ten level.
A simpler model that only uses his last two overall ratings predicts +1.29. Both of these would put
him a little below his 2023--24 level of +1.45. A very good player, but not quite the one from last
season.

\begin{table}[H]
\centering\small
\caption{A look at Leonard's season, part by part. Each value shows how far above (+) or below
($-$) an average NBA minute he was, where +1 is clearly good and +2 is elite. Reliability (0 to 1)
is how well each part carries over from one season to the next for players 31 and older (1,009
player seasons). The forecast predicts each part from its last two seasons.}
\begin{tabular}{lcccc}
\toprule
Part of his game & 2024--25 & 2025--26 & Reliability & 2026--27 forecast \\
\midrule
Scoring volume & $+1.47$ & $+2.83$ & 0.91 & $+2.14$ \\
Offensive load & $+0.93$ & $+1.85$ & 0.93 & $+1.46$ \\
Creation & $+0.17$ & $+0.81$ & 0.93 & $+0.49$ \\
Efficiency under load & $+0.33$ & $\mathbf{+1.48}$ & \textbf{0.72} & $+0.69$ \\
Box Plus/Minus & $+1.04$ & $+2.62$ & 0.82 & $+1.55$ \\
Win shares per 48 & $+0.67$ & $+1.87$ & 0.75 & $+1.01$ \\
Steals and blocks & $+0.52$ & $+0.78$ & 0.86 & $+0.52$ \\
Rebounding & $+0.11$ & $+0.35$ & 0.97 & $+0.21$ \\
Getting to the line & $-0.42$ & $+0.59$ & 0.86 & $+0.15$ \\
\midrule
\textbf{Impact score} & $+0.86$ & $\mathbf{+2.17}$ & & $\mathbf{+1.31}$ \\
\bottomrule
\end{tabular}
\label{tab:anatomy}
\end{table}

\section*{Can he stay in the top eight?}

The 644 has a 2-year blend, meaning that part of his 2025--26 season would carry forward.
Therefore, Kawhi doesn't have to repeat it to stay near the top. Over 65 games, he would need an
impact of +1.28 to be given a \gravity{} score of 69.3 (8th place score from last
edition).\src{5}{https://thegravityreport.com/notes/spec/} That bar is lower than you'd think: +1.28
is roughly the 12th-best per-minute season in the league, because his great 2025--26 still counts
toward next year's score. An impact of +1.1 would still keep him in the top 10. Our forecast of
+1.31 clears that top-eight bar by just 0.03, which is about as close as it gets. For players 31
and older, forecasts like this one usually miss by about 0.40 in either direction. Taking that into
account, his chance of staying top 8 is about 50\%, and top 10 about two in three.

In this situation, games matter less than they seem. The availability term gives last season's 65
games 30\% of the weight, so a 45-game season would actually lower the bar slightly, from +1.28 to
+1.25, because his great 2025--26 counts for more when next season is short. What's going to end up
moving the answer is whether the shooting holds or not.

There's also a simpler way to ask the question that doesn't depend on our model at all: will Kawhi
finish top eight in Box Plus/Minus (BPM), a public stat on Basketball Reference that estimates how
many points per 100 possessions a player adds over an average player? That version is harder for
him. It only counts next season, and he has to play at least 1,500 minutes to qualify. Our forecast
works out to a BPM of about 4.7. In the last eleven seasons, eighth place has never taken less than
5.1, and it has taken as much as 7.4. He's also only reached 1,500 minutes in six of his last eight
seasons.\src{6}{https://www.basketball-reference.com/leagues/NBA_2026_advanced.html}

\section*{Toronto}

Kawhi's trade finally got completed on September 14th after the investigation of his endorsement
deal closed.\src{7}{https://www.nba.com/news/raptors-clippers-kawhi-leonard-trade} The Raptors
extended him for 2 years and \$115
million.\src{8}{https://www.tsn.ca/nba/article/raptors-confirm-contract-extension-for-kawhi/}
Toronto's other creator, Scottie Barnes, used 23.4\% of the team's possessions last season. If
Kawhi's load falls in Toronto, the part of his surge that history says should last is the part that
the new role would take away. That's one way that the trade could change our forecast, and it
would definitely change it for the worse.

\section*{The forecast}

\forecastscale
\forecastitem{50}{Kawhi Leonard ranks top 8 in The 644, July 2027 edition.}{\#1--8 = Win \quad \#9 or below = Miss. This is our own ranking; we keep one bet on it each article so the model
gets graded too.}
\forecastitem{30}{Kawhi Leonard finishes 2026--27 top 8 in Box Plus/Minus among players with 1,500
regular-season minutes.}{Basketball Reference, final regular-season table.}
\forecastitem{60}{Kawhi Leonard uses 30\% or more of Toronto's possessions in 2026--27.}{The part of his
surge we expect to hold. He used 33.5\% last season. Basketball Reference usage rate; void if he plays
under 1,000 minutes.}
\forecastitem{65}{Kawhi Leonard's true shooting falls below .610 in 2026--27.}{The part we expect to fade.
He shot .629 last season. Basketball Reference; void if he plays under 1,000 minutes.}

\medskip\noindent{\footnotesize\textcolor{steel}{Public-stat bets are graded on Basketball Reference's final regular-season
tables by April 30, 2027. Registered October 6, 2026. Scored by Brier score in the
Forecast Ledger.}}

\section*{Method and sources}
\small
\gravity-RS (\gravity{} Retrospective) follows the specification's retrospective equation,
1977--78 to 2025--26. Late surge: age 31 or older, 1,500+ minutes, prior season 1,000+, improvement
of 0.6 or more. The part-by-part forecast predicts each part from its current and prior value,
fitted on players 31 and older (1,009 cases). Rank and percentile translations use 2025--26
players with 1,000+ minutes (percentiles) and 1,500+ minutes (ranks). Code:
\texttt{papers/kawhi/surge.py}, \texttt{papers/common/gravity\_rs.py},
\texttt{papers/common/pipeline.py}.
\begin{enumerate}[label={[\arabic*]}]
\item The Gravity Report, The 644, July 2026 edition (CSV). \url{https://thegravityreport.com/the644/the644-2026-07.csv}
\item Basketball Reference, league advanced and per-100 tables, 1977--78 to 2025--26. \url{https://www.basketball-reference.com/leagues/}
\item Basketball Reference, Kawhi Leonard. \url{https://www.basketball-reference.com/players/l/leonaka01.html}
\item Cleaning the Glass, Kawhi Leonard (garbage time filtered; percentiles by position), accessed September 29, 2026. \url{https://cleaningtheglass.com/stats/player/2173}
\item The Gravity Report, The \gravity{} Specification, v3. \url{https://thegravityreport.com/notes/spec/}
\item Basketball Reference, 2025--26 advanced table. \url{https://www.basketball-reference.com/leagues/NBA_2026_advanced.html}
\item NBA.com, Raptors--Clippers Leonard trade. \url{https://www.nba.com/news/raptors-clippers-kawhi-leonard-trade}
\item TSN, Raptors confirm Leonard extension. \url{https://www.tsn.ca/nba/article/raptors-confirm-contract-extension-for-kawhi/}
\end{enumerate}

\end{document}
